\documentclass{resume}

\begin{document}

\resumeheader{Márcio Sobral}{Desenvolvedor de Software}{%
  \contactlink{mailto:me@marciosobral.com}{me@marciosobral.com}\contactsep
  \contactphone
  \contactlink{https://linkedin.com/in/marciosobral}{linkedin.com/in/marciosobral}\contactsep
  \contactlink{https://github.com/marciosobral}{github.com/marciosobral}\contactsep
  \contactlink{https://marciosobral.com}{marciosobral.com}%
}

\cvsection{Resumo}

Desenvolvedor full stack com foco em backend e integrações entre sistemas. Trabalho com TypeScript, Node.js, NestJS, React e React Native, e atualmente desenvolvo agentes de IA para vendas, cobrança e atendimento.

\cvsection{Experiência}

\entry{Desenvolvedor Pleno}{Tecto Tecnologia}{Remoto}{Dez 2025 - presente}
\entrysummary{Atuo na Vendedor IA, a plataforma de agentes de IA da Tecto, usada por mais de 30 empresas.}
\begin{bullets}
  \item Sou o principal contribuidor do backend (NestJS, GraphQL e PostgreSQL), que hoje atende cerca de mil conversas e 10 mil mensagens por dia, com respostas baseadas em RAG.
  \item Trabalhei em mais de 10 integrações com ERPs, CRMs, e-commerces e softwares de clínicas e escolas.
  \item Criei do zero o agente de cobrança, que busca as faturas vencidas no ERP do cliente, negocia o acordo e finaliza o pagamento sozinho. Em menos de sete meses, recuperou mais de R\$~10 milhões.
  \item Conduzi implantações da plataforma e ataquei as falhas que mais paravam a operação. Fluxos que quebravam toda semana passaram a rodar semanas sem incidente.
\end{bullets}

\entry{Desenvolvedor Pleno}{Opportuna Tecnologia}{Remoto}{Dez 2024 - Dez 2025}
\entrysummary{Ecossistema de bilheteria Izy Pass, com site e app de compra, painel administrativo-financeiro, app do produtor, app de validação na portaria e venda na maquininha de cartão.}
\begin{bullets}
  \item Liderei a evolução das interfaces de compra e validação, encurtando o caminho até o ingresso.
  \item Reestruturei código legado com padrões de projeto e reduzi em 30\% as falhas reportadas pelos usuários.
  \item Participei das decisões técnicas sobre eventos com horário de funcionamento na expansão para kartódromos.
\end{bullets}

\entry{Desenvolvedor Júnior}{Opportuna Tecnologia}{Brasília, DF}{Jun 2024 - Dez 2024}
\begin{bullets}
  \item Construí um sistema de prontuário para clínicas, com fichas e formulários que geram PDFs e termos de consentimento assinados eletronicamente. O preenchimento ficou 50\% mais rápido.
  \item Desenvolvi o site da Opportuna e sites de clientes em Next.js, com SEO pensado desde o início do projeto.
\end{bullets}

\cvsection{Habilidades}

\skill{LINGUAGENS}{TypeScript\sep JavaScript\sep Python\sep Java\sep SQL}
\skill{INTERFACES}{Figma\sep React\sep Next.js\sep Tailwind CSS\sep shadcn/ui}
\skill{APLICATIVOS}{React Native\sep Expo\sep EAS\sep Kotlin\sep Sentry}
\skill{SERVIDOR E DADOS}{Node.js\sep NestJS\sep GraphQL\sep PostgreSQL\sep Redis\sep Bull}
\skill{IA E ATENDIMENTO}{OpenAI API\sep RAG\sep LangChain\sep n8n\sep Meta Cloud API}
\skill{INFRAESTRUTURA}{Docker\sep Vercel\sep Cloudflare\sep Railway\sep Render\sep VPS}
\skill{SEO E MÉTRICAS}{Search Console\sep Google Analytics\sep Tag Manager}

\cvsection{Projetos open source}

\project{comms-crm-core}{https://github.com/marciosobral/comms-crm-core}{CRM para revendas de telecom, com TanStack Start, NestJS, Prisma e PostgreSQL.}
\project{crm-core}{https://github.com/marciosobral/crm-core}{CRM enxuto para times de vendas, com Effect, React e PostgreSQL.}

\cvsection{Formação}

\entry{Pós-graduação em Segurança Cibernética}{UniCEUB}{}{Fev 2022 - Out 2022}
\begin{bullets}
  \item Ethical Hacking, Desenvolvimento Seguro de Sistemas e Gestão da Segurança Cibernética
\end{bullets}

\entry{Bacharelado em Ciência da Computação}{UniCEUB}{}{Jan 2018 - Dez 2021}
\begin{bullets}
  \item Talento Acadêmico 2020 (Prêmio ASSESPRO/DF e SINDESEI/DF)
  \item Prêmio Projeto Destaque 2021 no TCC, categoria Completude
\end{bullets}

\cvsection{Certificados}

\certificate{ESL Level VI (500h)}{St. George International College, Toronto}{Nov 2023}
\certificate{Curso de Oratória}{SA Treinamentos}{Mai 2023}

\cvsection{Idiomas}

Português (nativo)\sep Inglês (B2)

\end{document}
